\documentclass[12pt,a4paper]{report}

\usepackage[utf-8]{inputenc}
\usepackage[english]{babel}
\usepackage{amsmath}
\usepackage{amssymb}
\usepackage{graphicx}
\usepackage{xcolor}
\usepackage{listings}
\usepackage{hyperref}
\usepackage{geometry}
\usepackage{fancyhdr}

\geometry{margin=1in}
\pagestyle{fancy}
\fancyhead[L]{Xylella Nonlocal System Simulation}
\fancyfoot[C]{\thepage}

\lstset{
    language=Python,
    basicstyle=\ttfamily\small,
    keywordstyle=\color{blue},
    commentstyle=\color{gray},
    stringstyle=\color{red},
    breaklines=true,
    showstringspaces=false,
    tabsize=4
}

\title{\textbf{Xylella Nonlocal System Simulation} \\ \large A Comprehensive Python Program Documentation}
\author{Mauro Garavello}
\date{\today}

\begin{document}

\maketitle

\tableofcontents

\newpage

\chapter{Introduction}

The Xylella project is a sophisticated Python-based simulation program designed to model the dynamics of a nonlocal system, particularly focused on the spread and dynamics of \textit{Xylella fastidiosa} in olive tree populations. The program implements numerical methods for solving structured population models with spatial and age-structured components.

\section{Project Overview}

The simulation employs:
\begin{itemize}
    \item \textbf{Finite Volume Methods}: Specifically the Lax-Friedrichs and upwind schemes for solving conservation laws
    \item \textbf{Dimensional Splitting}: Separates the evolution in age, spatial (x), and spatial (y) dimensions
    \item \textbf{Age-Structured Populations}: Models juvenile and adult population dynamics separately
    \item \textbf{Spatial Coupling}: Incorporates nonlocal kernels and weeds convolution
    \item \textbf{Parallel Processing}: Uses multiprocessing for computational efficiency in both simulation and visualization
\end{itemize}

\chapter{Program Architecture}

\section{Directory Structure}

\begin{lstlisting}
Python/
  main_xylella.py           # Main entry point and orchestrator
  src/
    SimulationState.py      # Data structures for population states
    evolution.py            # Time-stepping evolution algorithm
    finite_volume_methods.py # Numerical flux implementations
    kernel_meshesgrids.py   # Kernel mesh generation utilities
    kernels.py              # Kernel definition and creation
    meshes.py               # Spatial and age mesh utilities
    mortality.py            # Mortality parameter management
    params_scheme.py        # Simulation parameter dataclass
    plot_functions.py       # Visualization and plotting
    save_load.py            # I/O operations for state persistence
  simulations/              # Simulation configurations
    Aveiro_LF/
    Aveiro_upwind/
  tests/                    # Test configurations
    velocity_down_no_source_LF/
    velocity_down_no_source_upwind/
    velocity_left_no_source_LF/
    velocity_left_no_source_upwind/
\end{lstlisting}

\section{Core Design Principles}

The program follows a modular architecture with clear separation of concerns:

\begin{enumerate}
    \item \textbf{State Management}: Population distributions are stored in specialized dataclasses with augmented arrays for boundary conditions
    \item \textbf{Numerical Methods}: Finite volume implementations are isolated in dedicated modules
    \item \textbf{Configuration}: Each simulation configuration is parameter-driven through external Python parameter files
    \item \textbf{I/O Operations}: State snapshots are saved at specified time intervals for later analysis and visualization
\end{enumerate}

\chapter{Core Modules}

\section{Main Entry Point: \texttt{main\_xylella.py}}

The main program orchestrates the entire simulation workflow and serves as the interface for users.

\subsection{Command-Line Arguments}

\begin{lstlisting}[language=bash]
python main_xylella.py <dir_name> [options]

Arguments:
  dir_name                    Directory containing parameters.py
  
Options:
  -s, --simulation            Execute the simulation
  -p, --plot                  Generate plots from saved simulation data
  -pn, --processors_number    Number of processors for parallel tasks
  -m, --movie                 Create movies from plots
  -cb, --color_bar            Add color bar to contour plots
\end{lstlisting}

\subsection{Simulation Workflow}

The main program follows this workflow:

\begin{enumerate}
    \item \textbf{Parameter Loading}: Dynamically imports a \texttt{parameters.py} module from the specified directory
    \item \textbf{Mesh Creation}: Generates spatial meshes for $x_1, x_2$ coordinates and age structure
    \item \textbf{State Initialization}: Creates initial conditions for populations $u$ (juveniles), $w$ (adults), and $z$ (trees)
    \item \textbf{Evolution Loop}: Iteratively advances the simulation in time, saving snapshots at regular intervals
    \item \textbf{Mass Conservation Tracking}: Computes and logs $L^1$ norms (masses) and $L^\infty$ norms throughout the simulation
\end{enumerate}

\subsection{Key Variables}

\begin{itemize}
    \item \textbf{$u$}: Juvenile population distribution, 3D array with dimensions (age, $x_1$, $x_2$)
    \item \textbf{$w$}: Adult population distribution, 3D array with dimensions (age, $x_1$, $x_2$)
    \item \textbf{$z$}: Tree density distribution, 2D array with dimensions ($x_1$, $x_2$)
    \item \textbf{Mortality parameters}: $\delta_u$, $\delta_w$, $\delta_z$, $\lambda$ (spatially and age-dependent)
    \item \textbf{Kernels}: $\eta_\mu$ (weed kernel), $\nabla\eta$ (gradient kernel), $\eta_z$ (tree kernel)
\end{itemize}

\section{Simulation State: \texttt{SimulationState.py}}

This module defines the data structures that hold the simulation's state variables.

\subsection{State3D Dataclass}

Represents age-structured population distributions with ghost cells for boundary conditions:

\begin{lstlisting}[language=Python]
@dataclass
class State3D:
    augmented_density: NDArray       # With ghost cells
    density: NDArray                 # Core domain
    augmented_density_a: NDArray     # Augmented in age dim
    augmented_density_x: NDArray     # Augmented in x dim
    augmented_density_y: NDArray     # Augmented in y dim
    density_a_left: NDArray          # Shifted view
    density_x_right: NDArray
    density_x_left: NDArray
    density_y_right: NDArray
    density_y_left: NDArray
\end{lstlisting}

Ghost cells are used for implementing boundary conditions, particularly for juveniles (where $u(a=0, x_1, x_2) = \beta(t, x_1, x_2)$, the fertility term).

\subsection{State2D Dataclass}

Similar to State3D but for 2D spatial distributions (trees), with only $(x_1, x_2)$ dimensions.

\subsection{SimulationState Dataclass}

Bundles all three population distributions plus velocity fields:

\begin{lstlisting}[language=Python]
@dataclass
class SimulationState:
    u: State3D           # Juveniles
    w: State3D           # Adults
    z: State2D           # Trees
    u_velocity_x: State2D
    u_velocity_y: State2D
    w_velocity_x: State2D
    w_velocity_y: State2D
\end{lstlisting}

\section{Evolution Algorithm: \texttt{evolution.py}}

The \texttt{one\_step\_evolution} function advances the system by one time step.

\subsection{Algorithm Overview}

\begin{enumerate}
    \item \textbf{Compute Convolutions}: In parallel, compute nonlocal interactions:
    \begin{itemize}
        \item $z * \eta_z$ (tree-mediated interactions)
        \item $(u/\Delta a) * \nabla\eta_1$ and $(u/\Delta a) * \nabla\eta_2$ (juvenile spreading)
        \item $(w/\Delta a) * \nabla\eta_1$ and $(w/\Delta a) * \nabla\eta_2$ (adult spreading)
    \end{itemize}
    
    \item \textbf{Calculate Velocity Fields}: 
    \begin{itemize}
        \item $\xi_u = \text{velocity\_juveniles}(-\text{weeds} \cdot \text{conv}, -\text{weeds} \cdot \text{conv})$
        \item $\xi_w = \text{velocity\_adults}(-(\text{weeds} + z*\eta_z) \cdot \text{conv}, \ldots)$
    \end{itemize}
    
    \item \textbf{Compute Time Step}: CFL condition determines $\Delta t$:
    \[\Delta t = \text{CFL} \cdot \frac{\min(\Delta x_1, \Delta x_2, \Delta a_J, \Delta a_A)}{\alpha}\]
    where $\alpha = \max(|\xi_u|, |\xi_w|, 1)$
    
    \item \textbf{Update Populations}: Apply finite volume method:
    \[u^{n+1} = u^n + \Delta t \cdot \text{RHS}_u\]
    \[w^{n+1} = w^n + \Delta t \cdot \text{RHS}_w\]
    \[z^{n+1} = (1 - \Delta t \cdot \lambda \cdot \delta_z(w)) \cdot z^n\]
\end{enumerate}

\section{Finite Volume Methods: \texttt{finite\_volume\_methods.py}}

Implements two schemes for spatial discretization: Lax-Friedrichs and upwind.

\subsection{Lax-Friedrichs Scheme}

For the $x$-direction:

\begin{lstlisting}[language=Python]
def lax_friedrichs_for_x(augmented_state, augmented_velocity_x, 
                         x_mesh_size, alpha):
    first_term = (augmented_velocity_x[2:, :] * augmented_state[:, 2:, :] 
                  - augmented_velocity_x[:-2, :] * augmented_state[:, :-2, :])
    second_term = alpha * (augmented_state[:, 2:, :] 
                           - 2*augmented_state[:, 1:-1, :] 
                           + augmented_state[:, :-2, :])
    return (first_term - second_term) / (2*x_mesh_size)
\end{lstlisting}

The Lax-Friedrichs method adds artificial diffusion (via the $\alpha$ parameter) for stability:
\[\text{Flux} = \frac{1}{2\Delta x}[F(u_{i+1}) - F(u_{i-1}) - \alpha(u_{i+1} - 2u_i + u_{i-1})]\]

\subsection{Upwind Scheme}

Implements donor-cell flux-corrected transport:

\begin{lstlisting}[language=Python]
def upwind_for_x(augmented_state, augmented_velocity_x, x_mesh_size):
    first_term = (augmented_velocity_x[2:, :] * augmented_state[:, 2:, :] 
                  - augmented_velocity_x[:-2, :] * augmented_state[:, :-2, :])
    second_term = -0.5 * np.absolute(augmented_velocity_x[1:-1, :] 
                                     + augmented_velocity_x[2:, :]) \
                  * (augmented_state[:, 2:, :] - augmented_state[:, 1:-1, :])
    third_term = -0.5 * np.absolute(augmented_velocity_x[1:-1, :] 
                                    + augmented_velocity_x[:-2, :]) \
                 * (augmented_state[:, :-2, :] - augmented_state[:, 1:-1, :])
    return (first_term + second_term + third_term) / (2*x_mesh_size)
\end{lstlisting}

The upwind scheme uses velocity-dependent weights for better handling of steep gradients.

\subsection{Age Evolution}

Simple upwind method in age dimension:

\begin{lstlisting}[language=Python]
def upwind_for_age(augmented_state, age_mesh_size):
    return (augmented_state[1:, :, :] - augmented_state[:-1, :, :]) / age_mesh_size
\end{lstlisting}

\section{Mesh Generation: \texttt{meshes.py}}

Creates spatial and age grids for discretization.

\subsection{Spatial Meshes}

\begin{lstlisting}[language=Python]
@dataclass
class SpaceMeshes:
    x1: NDArray         # Mesh points in x1 direction
    x2: NDArray         # Mesh points in x2 direction
    dx1: float          # Cell size in x1
    dx2: float          # Cell size in x2
    N_x1: int           # Number of grid points
    N_x2: int
\end{lstlisting}

\subsection{Age Meshes}

\begin{lstlisting}[language=Python]
@dataclass
class AgeMeshes:
    a_J: NDArray        # Juvenile age grid
    a_A: NDArray        # Adult age grid
    da_J: float         # Juvenile age cell size
    da_A: float         # Adult age cell size
    N_a_J: int          # Number of juvenile age classes
    N_a_A: int          # Number of adult age classes
\end{lstlisting}

\section{Kernel Definition: \texttt{kernels.py}}

Defines the nonlocal kernels used in convolutions.

\subsection{Kernel Functions}

The kernels represent spatial interactions:

\begin{itemize}
    \item $\eta_\mu(x_1, x_2)$: Weed convolution kernel
    \item $\nabla\eta(x_1, x_2)$: Gradient kernel for movement coupling
    \item $\eta_z(x_1, x_2)$: Tree interaction kernel
\end{itemize}

\begin{lstlisting}[language=Python]
@dataclass
class Kernels:
    eta_z: NDArray              # Tree kernel
    grad_eta: list              # [grad_eta_x1, grad_eta_x2]
\end{lstlisting}

\section{Visualization: \texttt{plot\_functions.py}}

Generates comprehensive visualizations of simulation results.

\subsection{Plotting Functions}

\begin{enumerate}
    \item \textbf{plot\_contour}: 2D contour plots with customizable color levels
    \item \textbf{plot\_surface}: 3D surface plots for population distributions
    \item \textbf{plot\_mass}: Line plots of mass evolution over time
    \item \textbf{plot\_density}: Master function generating all plot types for a time snapshot
    \item \textbf{movie}: Creates AVI movies from sequences of plots using mencoder
\end{enumerate}

\subsection{Parallel Plotting}

The program uses multiprocessing for parallel plot generation:

\begin{lstlisting}[language=Python]
for j in range(len(files_list)):
    elem = files_list[j]
    if (j >= procs):
        pp[j % procs].join()
    pp[j % procs] = Process(target=plot.plot_density, args=...)
    pp[j % procs].start()
\end{lstlisting}

\section{Data I/O: \texttt{save\_load.py}}

Handles persistence of simulation state using NumPy's compressed format.

\begin{lstlisting}[language=Python]
def save_status(dir_simulation, file_name, u, w, z, time):
    # Saves densities and current time to NPZ file
    
def load_status(file_name):
    # Loads densities and time from NPZ file
    
def save_mass(dir_simulation, times, mass_u, mass_w, mass_z, 
              l_infty_u, l_infty_w, l_infty_z):
    # Saves conservation metrics
\end{lstlisting}

\chapter{Numerical Methods}

\section{Mathematical Framework}

The program solves a system of partial differential equations with nonlocal terms:

\begin{equation}
\frac{\partial u}{\partial t} + \frac{\partial u}{\partial a} + \nabla \cdot (\xi_u u) = -\delta_u u
\end{equation}

\begin{equation}
\frac{\partial w}{\partial t} + \frac{\partial w}{\partial a} + \nabla \cdot (\xi_w w) = -\delta_w w
\end{equation}

\begin{equation}
\frac{\partial z}{\partial t} = -\lambda \delta_z(w_1) z
\end{equation}

where velocities are defined through nonlocal interactions:

\begin{equation}
\xi_u = v_u(-\text{weeds} \cdot \nabla\eta * (u/\Delta a), -\text{weeds} \cdot \nabla\eta * (u/\Delta a))
\end{equation}

\section{Dimensional Splitting}

The program uses dimensional splitting to separate the multi-dimensional problem into 1D subproblems:

\begin{enumerate}
    \item Age evolution: $\partial_a u = -\delta u$
    \item $x_1$ evolution: $\partial_t u + \partial_{x_1}(\xi_x u) = 0$
    \item $x_2$ evolution: $\partial_t u + \partial_{x_2}(\xi_y u) = 0$
\end{enumerate}

This approach reduces computational complexity and allows for flexible numerical scheme selection per dimension.

\section{CFL Stability Condition}

The Courant-Friedrichs-Lewy (CFL) condition ensures stability:

\[\Delta t \leq \text{CFL} \cdot \frac{\min(\Delta x_1, \Delta x_2, \Delta a)}{\max(|\xi|, 1)}\]

The code dynamically adjusts $\Delta t$ to ensure:
\begin{itemize}
    \item CFL condition is satisfied
    \item Time steps align with age grid (when possible)
    \item Final time is reached exactly
\end{itemize}

\chapter{Configuration and Parameters}

\section{Parameter File Structure}

Each simulation requires a \texttt{parameters.py} file in the simulation directory containing:

\begin{lstlisting}[language=Python]
# Domain specification
x1_limits = [x1_min, x1_max]
x2_limits = [x2_min, x2_max]
a_limits = [a_min, a_max]

# Grid resolution
N_x1 = 100      # Number of x1 nodes
N_x2 = 100      # Number of x2 nodes
N_a = 50        # Number of age nodes

# Temporal parameters
T = 10.0        # Final simulation time
cfl = 0.5       # CFL parameter
N_pictures = 20 # Snapshot intervals
LF = True       # Use Lax-Friedrichs (False for upwind)

# Initial conditions (functions)
def u0(a, x1, x2): ...   # Initial juvenile distribution
def w0(a, x1, x2): ...   # Initial adult distribution
def z0(x1, x2): ...      # Initial tree distribution

# Mortality rates
def delta_u(a, x1, x2): ...
def delta_w(a, x1, x2): ...
def delta_z(x1, x2, w): ...

# Boundary/source terms
def beta(t, x1, x2): ...      # Juvenile production
def s(t, x1, x2): ...         # New tree source
def lmbda(x1, x2): ...        # Tree loss parameter

# Kernels
def eta_mu(x1, x2): ...       # Weed kernel
def grad_eta(x1, x2): ...     # Gradient kernel
def eta_z(x1, x2): ...        # Tree kernel

# Velocity functions
def velocity_juveniles(conv_x, conv_y): ...
def velocity_adults(conv_x, conv_y): ...

# Convolution method (optional)
convolution = signal.oaconvolve  # or signal.convolve2d

# Plotting ranges (optional)
density_limits = [[u_min, u_max], [w_min, w_max], [z_min, z_max]]
rcount, ccount = 50, 50  # Surface plot resolution
\end{lstlisting}

\chapter{Usage Examples}

\section{Running a Simulation}

\begin{lstlisting}[language=bash]
# Execute simulation with 4 processors
python main_xylella.py ./simulations/Aveiro_LF -s -pn 4

# This will:
# - Load parameters from Aveiro_LF/parameters.py
# - Run the simulation
# - Save snapshots to Aveiro_LF/Simulation/
# - Generate log file Aveiro_LF/Simulation/log_Simulation.txt
\end{lstlisting}

\section{Plotting Simulation Results}

\begin{lstlisting}[language=bash]
# Generate plots from saved simulation data
python main_xylella.py ./simulations/Aveiro_LF -p -pn 4 -cb

# This will:
# - Load saved snapshots
# - Generate contour and surface plots
# - Add color bars to plots
# - Generate log file Aveiro_LF/Simulation/log_Plots.txt
\end{lstlisting}

\section{Creating Movies}

\begin{lstlisting}[language=bash]
# Create animation videos from plot sequences
python main_xylella.py ./simulations/Aveiro_LF -m

# This requires mencoder to be installed
# Creates AVI files for surface and contour plots
\end{lstlisting}

\chapter{Output Files}

\section{Simulation Output Structure}

\begin{lstlisting}
simulations/Aveiro_LF/Simulation/
  saving_000.npz              # Snapshot at t=0
  saving_001.npz              # Snapshot at t=T/19
  ...
  saving_019.npz              # Snapshot at t=T
  Mass.npz                    # All conservation data
  log_Simulation.txt          # Simulation log
  log_Plots.txt               # Plotting log (if -p used)
  log_Movie.txt               # Movie creation log (if -m used)
  plot_contour_juveniles_*.png
  plot_contour_adults_*.png
  plot_contour_trees_*.png
  plot_surface_juveniles_*.png
  plot_surface_adults_*.png
  plot_surface_trees_*.png
  plot_mass_*.png
  movie_*.avi                 # Generated movies
\end{lstlisting}

\section{Console Logging}

Key logged information includes:

\begin{itemize}
    \item Initialization timestamps
    \item Initial and final mass for each population
    \item Initial and final $L^\infty$ norms
    \item Checkpoint messages during simulation
    \item Processor count and timing information
    \item Total simulation duration
\end{itemize}

\chapter{Performance Considerations}

\section{Computational Complexity}

\subsection{Per Time Step}

\begin{itemize}
    \item \textbf{Convolutions}: $O(N^2 \log N)$ via FFT (using scipy.signal.oaconvolve)
    \item \textbf{Finite Volume Updates}: $O(N^3)$ for 3D arrays (age + 2D space)
    \item \textbf{Parallel Convolutions}: 3 independent FFT operations executed in parallel
\end{itemize}

\subsection{Memory Requirements}

For a typical configuration with $N_x = 100$, $N_a = 50$:

\begin{itemize}
    \item State3D: $\approx$ 100 MB (including augmented arrays and shifted views)
    \item State2D: $\approx$ 5 MB
    \item Kernels: $\approx$ 10 MB
    \item Total: $\approx$ 400-500 MB for simulation state
\end{itemize}

\section{Optimization Features}

\begin{enumerate}
    \item \textbf{Multiprocessing}: Convolutions computed in parallel (lines 48-50 of evolution.py)
    \item \textbf{Memory Views}: Shifted array views avoid data copying
    \item \textbf{Efficient Convolution}: Uses scipy's overlap-add method (oaconvolve)
    \item \textbf{Parallel Plotting}: Independent plot generations executed in parallel
\end{enumerate}

\chapter{Extension Points}

The modular design allows for easy extensions:

\section{Custom Velocity Functions}

Parameters can define custom velocity functions:

\begin{lstlisting}[language=Python]
def velocity_juveniles(conv_x, conv_y):
    return (custom_vx(conv_x), custom_vy(conv_y))
\end{lstlisting}

\section{Alternative Kernels}

Replace kernel definitions in parameters.py for different interaction models.

\section{Additional Mortality Terms}

Add age or space-dependent mortality functions directly in parameters.

\section{Custom Initial Conditions}

Define arbitrary initial distributions through parameter functions.

\section{Alternative Numerical Schemes}

Add new methods to \texttt{finite\_volume\_methods.py} and toggle via \texttt{LF} parameter.

\chapter{Debugging and Logging}

\section{Log Files}

Three main log files provide diagnostic information:

\begin{enumerate}
    \item \textbf{log\_Simulation.txt}: Numerical results and timing
    \item \textbf{log\_Plots.txt}: Visualization progress
    \item \textbf{log\_Movie.txt}: Movie generation status
\end{enumerate}

\section{Inspection Points}

Key variables to inspect when debugging:

\begin{itemize}
    \item Mass conservation (should remain relatively constant or decrease monotonically)
    \item $L^\infty$ norms (check for oscillations or blow-up)
    \item CFL time steps (should decrease if velocities increase)
    \item Convolution outputs (should be smooth without spurious oscillations)
\end{itemize}

\chapter{References and Mathematical Notes}

\section{Finite Volume Methods}

The program implements classical finite volume schemes for conservation laws:

\begin{itemize}
    \item Lax-Friedrichs: Robust, diffusive, unconditionally stable in the discrete sense
    \item Upwind: Preserves monotonicity, steeper gradients but less diffusion
\end{itemize}

\section{Nonlocal Models}

The use of convolution kernels implements nonlocal dispersal and interaction operators, commonly used in:

\begin{itemize}
    \item Spatial ecology models
    \item Epidemic models with long-distance transmission
    \item Population dynamics with aggregation effects
\end{itemize}

\section{Age-Structured Models}

The age dimension allows modeling ontogenetic changes in population dynamics, fundamental to:

\begin{itemize}
    \item Matrix population models
    \item McKendrick-Von Foerster equations
    \item Stage-structured population dynamics
\end{itemize}

\appendix

\chapter{Key Functions Reference}

\section{Main Control Flow}

\begin{tabular}{|l|l|}
\hline
\textbf{Function} & \textbf{Purpose} \\
\hline
\texttt{start\_log(file\_name)} & Initialize logging \\
\texttt{create\_directory(directory)} & Setup output directory \\
\texttt{one\_step\_evolution(...)} & Advance simulation by $\Delta t$ \\
\hline
\end{tabular}

\section{Mesh Functions}

\begin{tabular}{|l|l|}
\hline
\textbf{Function} & \textbf{Purpose} \\
\hline
\texttt{spatial\_mesh(x1\_limits, x2\_limits, N\_x1, N\_x2)} & Create spatial grid \\
\texttt{age\_mesh(a\_limits, N\_a)} & Create age grid \\
\texttt{meshgrid\_2D(x1, x2)} & Generate 2D meshgrid \\
\texttt{meshgrids\_3D(a, x1, x2)} & Generate 3D meshgrids \\
\hline
\end{tabular}

\section{Numerical Methods}

\begin{tabular}{|l|l|}
\hline
\textbf{Function} & \textbf{Purpose} \\
\hline
\texttt{lax\_friedrichs\_for\_x/y(...)} & Lax-Friedrichs flux \\
\texttt{upwind\_for\_x/y/age(...)} & Upwind flux \\
\hline
\end{tabular}

\end{document}
